\section{Introduction}
\label{sec:introduction}

The disappearance of distant objects behind a horizon is ordinarily described as
geometric occlusion by a curved surface.  Refraction provides a conceptually
different mechanism: a vertically varying refractive index bends rays and may
prevent rays emitted by the lower part of an object from reaching an observer,
even when the underlying boundary is flat.  Merely stating that refraction can
bend light, however, does not determine whether it can reproduce the specific
visibility pattern associated with a spherical surface.  That claim requires a
refractive-index profile and the ray family generated by it.

This article formulates that question as a concrete problem in geometrical
optics.  Consider a flat, opaque boundary with a smooth, vertically stratified
refractive index $n(h)$ above it.  We ask whether such a medium can reproduce
the principal horizon effects produced by a sphere of radius $R$: lower parts
of distant objects disappearing first, the hidden portion reappearing when the
observer is raised, and the corresponding dependence of the limiting visible
distance on observer and target height.

The question is deliberately separated into two parts.  The first is
mathematical: does any admissible profile $n(h)$ generate the required ray
geometry?  The second is physical: if such a profile exists, is its magnitude
and sign compatible with the refractive-index gradient of an ordinary
atmosphere?  Mathematical existence alone is therefore not treated as evidence
that the mechanism occurs under terrestrial conditions.

Within the paraxial approximation, a profile increasing exponentially with
height produces upward-curving rays and reproduces the leading-order spherical
horizon law.  The same exponential profile also admits an exact ray solution.
A conformal change of variables maps its two-dimensional optical metric to the
Euclidean exterior of a circle, explaining why a flat boundary can acquire the
same qualitative occlusion geometry as a curved one.  We then compare the exact
refractive and spherical horizon laws, examine local image geometry, and derive
a general reconstruction formula relating a prescribed horizon-distance law to
the stratified refractive index that realizes it.

The physical comparison gives the decisive distinction.  The globe-like
refractive construction requires $dn/dh>0$: the refractive index must increase
with altitude.  Under ordinary atmospheric conditions, refractivity instead
decreases with altitude, so that $dn/dh<0$.  Thus the desired optical occlusion
is mathematically constructible above a flat surface, but the required medium
has the opposite vertical trend from the normal atmosphere.

\subsection{Scope and assumptions}
\label{sec:scope}

The analysis uses geometrical optics and assumes a stationary, isotropic,
lossless medium whose refractive index depends only on height.  The lower
boundary is flat and perfectly opaque.  The profile $n(h)$ is taken to be
smooth in the region traversed by the rays, and ray propagation is reciprocal.
Diffraction, scattering, absorption, turbulence, and time-dependent atmospheric
structure are excluded from the model.  These assumptions isolate the question
of whether vertical refraction alone can reproduce horizon-like occlusion; they
are not intended as a complete model of atmospheric imaging.

The two-dimensional construction is developed first because every ray in a
vertically stratified medium lies in a vertical plane.  Three-dimensional image
geometry is considered separately when discussing transverse magnification and
possible distortion.  Throughout, $R$ denotes the radius of the spherical
surface whose horizon geometry is used as the comparison target, not an
assumption about the geometry of the flat optical model itself.
